% sections/body.tex — Лаб. 10: Подготовка прикладной библиотеки
\section*{Цель работы}

Получить практические навыки проектирования, реализации, сборки,
документирования и тестирования прикладной библиотеки на языке~C:
научиться выделять публичный интерфейс, скрывать детали реализации,
собирать разделяемую версию~\texttt{.so}, использовать её из~Python
через \texttt{ctypes}, вести историю разработки в~git, применять
Makefile, статические анализаторы и санитайзеры, а также подготавливать
автоматическую документацию средствами Doxygen.

\HRule

\section*{Вариант~4 — Частотный анализ текста}

Разработать библиотеку \textbf{freqlib} для частотного анализа
ASCII-текста.  Предусмотреть:
\begin{enumerate}
    \item подсчёт частот символов ASCII;
    \item подсчёт частот слов ограниченной длины;
    \item поиск top-$k$ наиболее частых слов.
\end{enumerate}

\HRule

\section*{1. Структура проекта}

\begin{lstlisting}[style=Bashstyle]
freqlib/
 include/
   freqlib.h         public API with Doxygen comments
 src/
   freqlib.c         library implementation
   main.c            demo application
 tests/
   test_freqlib.c    C unit tests  (43 cases)
   test_freqlib.py   Python ctypes tests  (22 cases)
 docs/
   Doxyfile          Doxygen configuration
 reports/            static-analysis output  (.gitignore)
 build/              compiled artifacts       (.gitignore)
 Makefile
 README.md
 .gitignore
\end{lstlisting}

\HRule

\section*{2. Публичный интерфейс библиотеки}

\subsection*{2.1 Коды ошибок}

\begin{center}
\begin{tabular}{ll}
\toprule
Константа & Значение \\
\midrule
\texttt{FREQ\_OK}~(0)       & Операция выполнена успешно \\
\texttt{FREQ\_ERR\_NULL}~(-1)  & Передан \texttt{NULL}-указатель \\
\texttt{FREQ\_ERR\_EMPTY}~(-2) & Пустая входная строка или нет слов \\
\texttt{FREQ\_ERR\_ALLOC}~(-3) & Ошибка выделения памяти \\
\texttt{FREQ\_ERR\_PARAM}~(-4) & Некорректное числовое значение \\
\bottomrule
\end{tabular}
\end{center}

\subsection*{2.2 Типы данных}

\begin{lstlisting}[style=Cstyle]
/* Character frequency table (stack-allocated, no heap) */
typedef struct {
    size_t counts[128]; /* counts[c] == occurrences of ASCII char c */
    size_t total;       /* total ASCII characters seen               */
} char_freq_t;

/* One word + its occurrence count */
typedef struct {
    char   word[64];    /* NUL-terminated, max 63 bytes              */
    size_t count;
} word_entry_t;

/* Dynamic word-frequency table (heap-allocated entries) */
typedef struct {
    word_entry_t *entries;
    size_t        size;   /* distinct words stored */
    size_t        cap;    /* allocated capacity    */
} word_freq_t;
\end{lstlisting}

\subsection*{2.3 Функции}

\begin{lstlisting}[style=Cstyle]
int  freq_char_count(const char *text, char_freq_t *out);
void freq_char_free (char_freq_t *cf);

int  freq_word_count(const char *text, word_freq_t *out);
void freq_word_free (word_freq_t *wf);

int  freq_top_k(const word_freq_t *wf, size_t k,
                word_entry_t **result, size_t *result_n);
\end{lstlisting}

\subsection*{2.4 Правила владения памятью}

\begin{itemize}
    \item \texttt{freq\_char\_count} заполняет структуру на стеке вызывающего;
          освобождение — \texttt{freq\_char\_free} (обнуляет поля).
    \item \texttt{freq\_word\_count} выделяет \texttt{wf->entries} через
          \texttt{malloc}/\texttt{realloc};
          освобождение — \texttt{freq\_word\_free}.
    \item \texttt{freq\_top\_k} выделяет новый массив результата;
          вызывающая сторона освобождает его через \texttt{free()}.
\end{itemize}

\HRule

\section*{3. Описание алгоритмов}

\textbf{Подсчёт символов.}
За один проход строки инкрементируется \texttt{counts[c]} для каждого
символа~\texttt{c} со значением меньше~128.
Сложность: $O(n)$, где $n$~--- длина строки.

\textbf{Подсчёт слов.}
Два указателя: пропуск не-алфавитно-цифровых символов,
затем сбор «слова» (до~63~байт).
Слово ищется в динамическом массиве линейным поиском;
при совпадении счётчик инкрементируется,
при отсутствии добавляется новая запись
(ёмкость удваивается начиная с~16).
Сложность: $O(m \cdot n)$ в худшем случае,
где $m$~--- число слов, $n$~--- число различных слов.

\textbf{Top-k.}
Создаётся копия массива \texttt{entries}, сортируется по убыванию
счётчика через \texttt{qsort} (при равенстве~--- по алфавиту),
возвращаются первые $\min(k,\,|\textit{wf}|)$ элементов.
Сложность: $O(n \log n)$.

\HRule

\section*{4. Фрагменты ключевых файлов}

\subsection*{4.1 Заголовочный файл \texttt{include/freqlib.h} (фрагмент)}

\begin{lstlisting}[style=Cstyle]
#ifndef FREQLIB_H
#define FREQLIB_H

#include <stddef.h>

#define FREQ_OK         0
#define FREQ_ERR_NULL  -1
#define FREQ_ERR_EMPTY -2
#define FREQ_ERR_ALLOC -3
#define FREQ_ERR_PARAM -4

#define FREQ_ASCII_SIZE  128
#define FREQ_WORD_MAXLEN  63

typedef struct {
    size_t counts[128];
    size_t total;
} char_freq_t;

int  freq_char_count(const char *text, char_freq_t *out);
void freq_char_free (char_freq_t *cf);

/* word_freq_t, word_entry_t, freq_word_count,
   freq_word_free, freq_top_k  ... */
#endif
\end{lstlisting}

\subsection*{4.2 Реализация \texttt{freq\_word\_count} (ключевой фрагмент)}

\begin{lstlisting}[style=Cstyle]
const char *p = text;
while (*p != '\0') {
    while (*p != '\0' && !isalnum((unsigned char)*p)) p++;
    if (*p == '\0') break;

    char buf[FREQ_WORD_MAXLEN + 1];
    size_t len = 0;
    while (*p != '\0' && isalnum((unsigned char)*p)) {
        if (len < FREQ_WORD_MAXLEN) buf[len++] = *p;
        p++;
    }
    buf[len] = '\0';

    int idx = wf_find(out, buf);
    if (idx >= 0) {
        out->entries[idx].count++;
    } else {
        if (out->size == out->cap) wf_grow(out);
        memcpy(out->entries[out->size].word, buf, len + 1);
        out->entries[out->size].count = 1;
        out->size++;
    }
}
\end{lstlisting}

\subsection*{4.3 Makefile (фрагмент)}

\begin{lstlisting}[style=Bashstyle]
CC     := gcc
CFLAGS := -Wall -Wextra -Wpedantic -std=c11 -Iinclude

shared:
    $(CC) $(CFLAGS) -O2 -fPIC -shared -o build/libfreqlib.so src/freqlib.c

sanitize:
    $(CC) $(CFLAGS) -g -fsanitize=address,undefined \
          -o build/test_san tests/test_freqlib.c src/freqlib.c -Iinclude
    build/test_san
\end{lstlisting}

\subsection*{4.4 Python-вызов через \texttt{ctypes} (фрагмент)}

\begin{lstlisting}[style=Pystyle]
import ctypes
lib = ctypes.CDLL("build/libfreqlib.so")

class CharFreq(ctypes.Structure):
    _fields_ = [("counts", ctypes.c_size_t * 128),
                ("total",  ctypes.c_size_t)]

lib.freq_char_count.argtypes = [ctypes.c_char_p,
                                  ctypes.POINTER(CharFreq)]
lib.freq_char_count.restype  = ctypes.c_int

cf = CharFreq()
rc = lib.freq_char_count(b"hello world", ctypes.byref(cf))
assert rc == 0
assert cf.counts[ord('l')] == 3
lib.freq_char_free(ctypes.byref(cf))
\end{lstlisting}

\HRule

\section*{5. Команды сборки и результаты}

\begin{lstlisting}[style=Bashstyle]
$ make all
[shared]  build/libfreqlib.so
[app]     build/app
43 passed, 0 failed
Build complete.

$ make run
freqlib demo
Character frequencies (printable ASCII, count > 0):
  ' '    67
  'e'    31   't'    28   'o'    27   's'    23
  'a'    22   ...
  total: 331
Distinct words: 48
Top 10 most frequent words:
   1.  to                    7
   2.  a                     3
   3.  and                   3
   4.  the                   3
   5.  The                   2
   ...
Done.

$ make syntax
[syntax]  OK - no errors

$ make sanitize
43 passed, 0 failed   (AddressSanitizer + UBSan: clean)

$ make py-test
=== freqlib Python test suite ===
22 passed, 0 failed

$ make analyze
--- gcc -fanalyzer ---   (no warnings)
--- cppcheck ---         (no errors)

$ make docs-html
[docs]    docs/html/index.html
\end{lstlisting}

\HRule

\section*{6. Таблица тестовых случаев}

\begin{center}
\renewcommand{\arraystretch}{1.12}
\small
\begin{longtable}{p{4.2cm}|p{3.5cm}|p{3.2cm}|c}
\toprule
Входные данные & Функция & Ожидаемый результат & Итог \\
\midrule
\texttt{NULL} вместо текста      & \texttt{freq\_char\_count} & \texttt{FREQ\_ERR\_NULL}  & PASS \\
\texttt{NULL} вместо \texttt{out}& \texttt{freq\_char\_count} & \texttt{FREQ\_ERR\_NULL}  & PASS \\
Пустая строка \texttt{""}        & \texttt{freq\_char\_count} & \texttt{FREQ\_ERR\_EMPTY} & PASS \\
\texttt{"aabbc"}                 & \texttt{freq\_char\_count} & counts: a=2, b=2, c=1    & PASS \\
Один символ \texttt{"z"}         & \texttt{freq\_char\_count} & counts[z]=1, total=1      & PASS \\
Три пробела \texttt{"   "}       & \texttt{freq\_char\_count} & counts[' ']=3, total=3    & PASS \\
После \texttt{freq\_char\_free}  &                            & total = 0                 & PASS \\
\texttt{NULL} вместо текста      & \texttt{freq\_word\_count} & \texttt{FREQ\_ERR\_NULL}  & PASS \\
Пустая строка                    & \texttt{freq\_word\_count} & \texttt{FREQ\_ERR\_EMPTY} & PASS \\
Только знаки \texttt{"!!! ???"}  & \texttt{freq\_word\_count} & \texttt{FREQ\_ERR\_EMPTY} & PASS \\
\texttt{"the cat and the dog"}   & \texttt{freq\_word\_count} & size=4, the:2             & PASS \\
Одно слово \texttt{"hello"}      & \texttt{freq\_word\_count} & size=1, count=1           & PASS \\
Знаки препинания                 & \texttt{freq\_word\_count} & hello:2                   & PASS \\
Слово 100 символов               & \texttt{freq\_word\_count} & обрезается до 63, OK      & PASS \\
Три регистра одного слова        & \texttt{freq\_word\_count} & 3 разные записи           & PASS \\
\texttt{k=0}                     & \texttt{freq\_top\_k}      & \texttt{FREQ\_ERR\_PARAM} & PASS \\
\texttt{NULL} аргумент           & \texttt{freq\_top\_k}      & \texttt{FREQ\_ERR\_NULL}  & PASS \\
\texttt{"a b a c a b"}, k=2      & \texttt{freq\_top\_k}      & [a:3, b:2]                & PASS \\
$k >$ числа слов                 & \texttt{freq\_top\_k}      & все слова                 & PASS \\
$k=1$                            & \texttt{freq\_top\_k}      & только «a»                & PASS \\
Пустая таблица                   & \texttt{freq\_top\_k}      & resn=0, res=NULL          & PASS \\
\bottomrule
\multicolumn{4}{r}{\small Итого: \textbf{43} C-теста, \textbf{22} Python-теста --- все пройдены} \\
\end{longtable}
\end{center}

\HRule

\section*{7. История разработки (\texttt{git log --oneline})}

\begin{lstlisting}[style=Bashstyle]
8d2e082 docs: add changelog; note fix for no-word input edge case
09b7b0c test(python): add ctypes integration tests with explicit argtypes/restype
5c01477 test(c): add unit tests covering typical, boundary and error cases
be225f6 build: add demo application and Makefile with all targets
d504960 feat(src): implement freq_char_count, freq_word_count and freq_top_k
e2aa5ae feat(include): add public header with types, error codes and Doxygen comments
67cc512 init: scaffold project structure, .gitignore and README
\end{lstlisting}

\HRule

\section*{8. Выводы}

В ходе выполнения лабораторной работы~№\,10 реализована библиотека
\textbf{freqlib} для частотного анализа ASCII-текста (вариант~4).

\medskip
\noindent\textbf{Реализовано:}
\begin{itemize}
    \item Публичный интерфейс в \texttt{include/freqlib.h}
          с полными Doxygen-комментариями и единой схемой кодов ошибок.
    \item Три функциональных группы API:
          \texttt{freq\_char\_count}, \texttt{freq\_word\_count},
          \texttt{freq\_top\_k}.
    \item Разделяемая библиотека \texttt{libfreqlib.so}
          (собирается с \texttt{-fPIC -shared}).
    \item Демонстрационное приложение \texttt{src/main.c}.
    \item \textbf{43} C-теста и \textbf{22} Python-теста через
          \texttt{ctypes} --- все пройдены.
    \item Makefile с~12~целями:
          \texttt{all, shared, app, run, test, py-test,}
          \texttt{syntax, analyze, sanitize, docs-html, docs-pdf, clean}.
    \item История разработки в git: 7~осмысленных коммитов.
\end{itemize}

\medskip
\noindent\textbf{Найденные ошибки и исправления:}
\begin{itemize}
    \item Изначально \texttt{freq\_word\_count} возвращала \texttt{FREQ\_OK}
          для строк, не содержащих слов (например, \texttt{"!!! ???"}); 
          исправлено --- теперь возвращается \texttt{FREQ\_ERR\_EMPTY}.
    \item Статический анализ (\texttt{cppcheck}, \texttt{gcc~-fanalyzer})
          и санитайзеры (ASan~+~UBSan) ошибок не выявили.
\end{itemize}

\medskip
\noindent\textbf{Ограничения и возможные расширения:}
\begin{itemize}
    \item Поиск слова линейный $O(n)$ --- при больших текстах
          можно заменить хеш-таблицей.
    \item Сравнение слов чувствительно к регистру ---
          можно добавить нечувствительный режим.
    \item Максимальная длина слова фиксирована (63~байта) ---
          можно сделать настраиваемой через параметр функции.
\end{itemize}
